\subsection{Tóm lược (Summary)}
Các kỹ thuật DFT chuyển đổi mạch từ cấu trúc hộp đen hoặc mô hình trạng thái (State Machine, Huffman Model) phức tạp thành luồng mạch tổ hợp thuần túy để thực thi việc sinh mẫu kiểm tra (ATP Generation). 

Trong thực tế, quy trình chèn quét scan được vận hành theo các bước: 
\begin{enumerate}
    \item Chạy sơ đồ chức năng (RTL);
    \item Nhanh chóng đưa chuỗi (Partial hoặc Full Scan) qua công cụ tổng hợp;
    \item Cắt gãy vòng phản hồi (Break Feedback Loop) thông qua Unfolding;
    \item Nạp tổ hợp MUX-DFF (Dùng NbarT chuyển chế độ);
    \item Tạo Testbench hoặc bộ kiểm thử Virtual Tester bằng Verilog;
    \item Nén và chạy tối giản chuỗi Test Vector thông qua các kĩ thuật Test Compaction tiên tiến nhất. Tối ưu thời gian là bài toán chi phí (Tradeoff) giữa số chân I/O vật lý trên bo mạch (Package pins) và số Clock Cycle thực tế mỗi khi kích hoạt kiểm duyệt chip.
\end{enumerate}

\subsection{Hướng phát triển tương lai (Future Work)}
Mặc dù Scan Design đã định hình tiêu chuẩn công nghiệp (De Facto) cho quá trình kiểm thử vi mạch, sự phát triển không ngừng của định luật Moore và thiết kế SoC (System-on-Chip) đang thúc đẩy lĩnh vực DFT tiến tới các công nghệ nén và tự động hóa mức độ cực cao:
\begin{itemize}
    \item \textbf{Nén dữ liệu kiểm thử (Test Compression):} Khi số lượng Flip-Flop trong hệ thống chạm ngưỡng hàng chục triệu, kỹ thuật Multiple Scan cũng đối mặt với "nút thắt cổ chai" băng thông từ máy ATE. Các cấu trúc giải nén (Decompressor) và nén phản hồi (Compactor) được chế tạo sẵn trên chip (On-chip) giúp bơm dữ liệu kiểm thử khổng lồ qua số lượng kênh giới hạn, giảm 10-100 lần bộ nhớ ATE và Test Time.
    \item \textbf{Tự kiểm thử tích hợp (Built-In Self-Test - BIST):} Loại bỏ hoàn toàn sự phụ thuộc vào các hệ thống máy ATE đắt đỏ bên ngoài. BIST nhúng bộ sinh vector tự động (PRPG) và phân tích lỗi (MISR) vào tận sâu bên trong bộ vi xử lý. Hệ thống có khả năng tự chẩn đoán bệnh ngay trên bo mạch (At-speed testing), phục vụ kiểm định thời gian thực cho hạ tầng mạng 5G/6G, xe tự lái (Automotive) hay hàng không vũ trụ.
\end{itemize}
